\begin{abstract}
\label{sec:abstract}

Detecting that a scientific-agent trajectory is failing does not determine
whether evidence is worth purchasing or which recovery action is safe. This
paper introduces a cost-sensitive decision framework that, given calibrated
failure and cause beliefs from a frozen measurement benchmark, selects
verification and recovery actions under an explicit utility-and-cost budget.
We formalize the sequential decision problem over a partially observed
trajectory state, define an action space spanning verify, continue, recover
(retry, repair, rollback, re-plan), escalate, abstain, and terminate, and
specify action-specific costs and postcondition checks enforced by an
independent adjudicator. A randomized, matched-budget online study compares
six policies—never verify, always verify, fixed periodic, random, threshold,
and adaptive—on the primary estimand of independently verified scientific
success at matched cost. We then construct a verified-intervention dataset
from independently adjudicated transitions and test whether
process-supervised training on those transitions improves unseen-fault
recovery and selective risk over recovery SFT and outcome-only
post-training, under a clean-task non-inferiority gate. The controller, the
intervention dataset, and the training evaluation are built on a frozen
benchmark contract, and all reward, verifier, and evaluation channels are
audited for independence.
\end{abstract}
